\section{Does the Story Continue? A Covariant Unification in 12d}
\label{sec:covariant-unification}

In the previous section we saw that the D(-1) and D7 solutions lift to 12d, and that, if a 12d theory existed, reducing its covariantly constant spinor condition on a torus would reproduce the type IIB Killing spinor equations.
It is natural to ask whether other brane-like solutions of type IIB lift in the same way.


\subsection{Brane-like solutions}
\label{sec:general brane solns}
Working in the Einstein frame, with $F_n$ an n-form and the action
\begin{equation}
S = \int d^Dx\sqrt{-g}\left[
    R - \frac{1}{2}(\partial\phi)^2 - \frac{1}{2}e^{\alpha\phi}|F_n|^2
\right],
\end{equation}
the equations of motion, together with the Bianchi identity (assuming $F = dA$), are given by \cite{Stelle:1998xg}
\begin{equation}
\begin{aligned}
R_{MN}
&= \frac{1}{2}\partial_M\phi\partial_N\phi + e^{\alpha\phi}S_{MN}\big|_F,\\
dF&=0,\\
d[e^{\alpha\phi}*_D F] &=0,\\
\Box\phi&=\frac{\alpha}{2}e^{\alpha\phi}|F|^2,
\end{aligned}
\end{equation}
where
\begin{equation}
S_{MN}\big|_{F}\equiv
\frac{1}{2}\left[
    \frac{1}{(n-1)!}F_{M...}F_N^{...}
    -\frac{n-1}{D-2}\left(\frac{1}{n!}F^2\right)g_{MN}
\right],
\end{equation}
and $\Delta\equiv \alpha^2 +\frac{2 d \tilde{d}}{D-2}$, with $\tilde{d} \equiv D-d-2$.
There exists the electric solution, defined as the case where the form-field potential has legs along the worldvolume directions,
\begin{equation}
\begin{aligned}
ds_D^2
&=H^{-\frac{4\tilde d}{\Delta(D-2)}}ds_{d}^2+
H^{\frac{4d}{\Delta(D-2)}}ds_{D-d}^2,\quad
d = n-1,\\
F_{\mu_0\mu_1...\mu_{n-2},m}
&=\frac{2}{\sqrt\Delta}\varepsilon_{\mu_0\mu_1...\mu_{n-2}}\partial_m (H^{-1}),\\
e^\phi
&=H^{2\alpha/\Delta},
\end{aligned}
\end{equation}
as well as the magnetic solution, defined as the case where the field strength has its legs along the transverse directions,
\begin{equation}
\begin{aligned}
ds_D^2
&=H^{-\frac{4\tilde d}{\Delta(D-2)}}ds_{d}^2+
H^{\frac{4d}{\Delta(D-2)}}ds_{D-d}^2,\quad
d = D-n-1,\\
F_{m_1...m_n}
&= \frac{2}{\sqrt{\Delta}}\varepsilon_{m_1...m_n m}\partial_m H,\\
e^\phi &= H^{-2\alpha/\Delta}.
\end{aligned}
\end{equation}
In 11d supergravity, type IIA, and type IIB, one finds $\Delta = 4$ for all branes.
We provide a catalog of the various brane solutions in Appendix~\ref{sec:catalog-of-branes}.

One can readily see that there do not exist nontrivial ways to embed the type IIB metric, scalars, and form fields into a 12d metric and form fields.
At the level of supergravity, 12d fields transform under the little group SO(10) while the 10d fields transform under SO(8).
In Table~\ref{tab:10d-12d-dof} we summarise the 10d massless on-shell degrees of freedom in type IIB, alongside the on-shell degrees of freedom of the corresponding 12d field of the same form rank.

\begin{table}[h]
\centering
\begin{tabular}{l|ccccccc}
\hline
 & metric & scalar & 1-form & 2-form & 3-form & 4-form & 5-form \\
\hline
Type IIB $SO(8)$ & $1\times 35$ & $2\times 1$ & $0\times 8$ & $2\times 28$ & $0\times 56$ & $1\times 35$ & $-$ \\
12d $SO(10)$ & $54$ & $-$ & $10$ & $45$ & $120$ & $210$ & $126$ \\
\hline
\end{tabular}
\caption{
On-shell bosonic degrees of freedom of the type IIB fields, given as (number of fields) $\times$ (dimension of the $SO(8)$ little-group representation), compared with the dimension of the $SO(10)$ representation of a 12d field of the same form rank.
}
\label{tab:10d-12d-dof}
\end{table}
As the table shows, the type IIB degrees of freedom do not fill out the SO(10) representations of the 12d fields. 
Type IIB therefore cannot arise as the dimensional reduction of a genuine 12d (super)gravity that realises the full 12d spacetime symmetries.

This does not, however, exhaust the 12d interpretation.
Already in the axio-dilaton sector the embedding into the 12d metric was only partial: the two off-diagonal Kaluza-Klein 1-forms were set to zero, so the full 12d metric was never reconstructed. It nonetheless provided a 12d perspective on the D(-1) and D7 branes, as we saw in the previous section.
In the same spirit, the following subsection organises the type IIB fields into 12d objects \emph{without} realising the full 12d spacetime symmetries: not all degrees of freedom of the 12d graviton and form fields are excited, so the construction is best understood as a covariant repackaging of the 10d theory rather than the reduction of a genuine 12d one.

\subsection{A Covariant Unification in 12d}
\label{sec:a-covariant-unification-in-12d}

The various brane-related evidence for an effective 12d interpretation of the type IIB theory suggests a very specific 12d interpretation of the axio-dilaton sector.
In addition, one seeks a 12d interpretation of the type IIB RR and NSNS form fields.
In this section, we gather the various 12d insights obtained in the previous section from the type IIB branes, and propose a covariant 12d unification of the $\SLR{}\times SO(9,1)$ form fields.
We will use hats $\widehat{\phantom{x}}$ to denote fields in the higher dimension.
The 10d metric $g_{mn}$ with $m,n=0,1,...,9$ will be interpreted as embedded inside a 12d metric $\widehat{g}_{MN}$, with $M,N=0,1,...,11$. We will use $\widehat R$ to denote the Ricci scalar computed with $\widehat{g}_{MN}$, and use $\widehat{F}_{n+1} = d\widehat{C}_{n}$ to denote form fields in 12d.

Let the 12d coordinates be parameterised by $(x^m,z_1,z_2)$, and let $\widehat{g}_{ab}$ be the $2\times 2$ metric on the torus.\footnote{
We orient the torus so that $dz_2\wedge dz_1$ is the positively-oriented volume element, i.e. $\int_{T_2}dz_1\wedge dz_2 = -\mathrm{Vol}(T_2)$. This fixes the sign of the 12d Hodge dual $*_{12}$, and hence the signs in the self-duality relation \eqref{12d interpretation of 10d self-duality} and in the Chern-Simons term.
}
The key insight from the previous section is that the axio-dilaton sector of type IIB may arise from the 12d metric embedding given by
\begin{equation}
\widehat{g}_{MN}=\begin{pmatrix}
g_{mn} & 0 \\
0 &\widehat{g}_{ab}
\end{pmatrix},\quad
\widehat{g}_{ab}
=\frac{1}{\tau_2}\begin{pmatrix}
1& \tau_1\\
\tau_1 & \tau_1^2+\tau_2^2
\end{pmatrix}.
\label{12d metric ansatz}
\end{equation}
In particular,
\begin{equation}
\begin{aligned}
\frac{1}{2\kappa^2}\int d^{10}x\sqrt{-g}\left(
R - \frac{|\nabla \tau|^2}{2\tau_2^2}
\right)
&=
\frac{\mathrm{Vol}(T_2)}{2\kappa_{12}^2}\int d^{10}x \sqrt{-\widehat{g}} \widehat{R}\\
&=\frac{1}{2\kappa_{12}^2}\int_{T_2}dz_1 dz_2\int d^{10}x\sqrt{-\widehat{g}}\widehat{R},
\end{aligned}
\label{axio-dilaton 12d uplift}
\end{equation}
where we have schematically defined $\kappa_{12}$ by 
\begin{equation}
\frac{1}{\kappa^2} = \frac{\mathrm{Vol}(T_2)}{\kappa_{12}^2}
=\frac{1}{\kappa_{12}^2}\int_{T_2}*_2 1,\quad
[\kappa_{12}^2] = L^{10}.
\end{equation}
The 3-form field strengths form an \SLR{} doublet.
To unify them in 12d we define a 12d 4-form field strength with exactly one leg on the torus:
\begin{equation}
\widehat{F}_4=d\widehat C_3,\quad
\widehat C_3 \equiv B_2\wedge dz_1 + C_2\wedge dz_2.
\end{equation}
Using the 12d metric ansatz \eqref{12d metric ansatz}, the 10d 3-form field strengths and their axio-dilaton couplings follow from contracting $\widehat F_4$:
\begin{equation}
\begin{aligned}
|\widehat{F}_4|^2\bigg|_{\widehat{g}_{MN}}
&=\widehat{g}^{z_1 z_1}|H_3|^2+2\widehat{g}^{z_1 z_2}(F_3\cdot H_3)+\widehat{g}^{z_2 z_2}|F_3|^2\\
&=\left(
e^{-\Phi}|H_3|^2 + e^\Phi|F_3-C_0H_3|^2
\right)\bigg|_{g_{mn}}.
\end{aligned}
\end{equation}
Finding a 12d interpretation for the 5-form sector is trickier. 
One observes that in 10d, the composite \SLR{} singlet 5-form 
\begin{equation}
\tilde F_5 = F_5 - \frac12 C_2 \wedge H_3 + \frac{1}{2}B_2\wedge F_3    
\end{equation}
cannot be sensibly constructed in 12d, due to the lack of a pair of form field potential and strength with ranks that sum to 5 in 12d.
However, the 10d five-form field strength admits two possible uplifts to 12d:
\begin{equation}
\widehat{F}_5 = F_5,\quad \widehat{F}_7=F_5\wedge dz_1 \wedge dz_2.
\end{equation}
Then note that
\begin{equation}
\begin{aligned}
|\widehat{F}_5|^2\bigg|_{\widehat{g}_{MN}} 
&= |F_5|^2\bigg|_{g_{mn}},\\
|\widehat{F}_7|^2\bigg|_{\widehat{g}_{MN}}
&=\det(\widehat{g}^{ab})|F_5|^2\bigg|_{g_{mn}} = |F_5|^2\bigg|_{g_{mn}}.
\end{aligned}
\end{equation}
Then we may define a 12d 7-form
\begin{equation}
\widetilde{\widehat{F}}_7 \equiv \widehat{F}_7 - \frac{1}{2}\widehat{C}_3 \wedge \widehat{F}_4
\end{equation}
that exactly supplies the type IIB composite 5-form contribution upon contraction in 12d:
\begin{equation}
|\widetilde{\widehat{F}}_7|^2\bigg|_{\widehat{g}_{MN}} =
\det(\widehat{g}^{ab})|\tilde F_5|^2\bigg|_{g_{mn}}.
\end{equation}
The additional utility of the composite seven-form is that it lets us write the genuine 10d self-duality condition on the composite five-form $\tilde F_5$ as a 12d Hodge duality condition
\begin{equation}
\widetilde{\widehat{F}}_7 = -*_{12}\widetilde{\widehat{F}}_5\quad
\Leftrightarrow\quad
\tilde F_5 = *_{10}\tilde F_5.
\label{12d interpretation of 10d self-duality}
\end{equation}
The 10d Chern-Simons term can also be obtained using the 12d potentials and their corresponding field strengths:
\begin{equation}
\frac{1}{2}\int_{T_2\times \mathbb{R}^{1,9}} \widehat{C}_4\wedge \widehat{F}_4\wedge\widehat{F}_4=
\mathrm{Vol}(T_2)\int_{\mathbb{R}^{1,9}}C_4\wedge H_3\wedge F_3.
\end{equation}
We note that reproducing the type IIB Chern-Simons term in 12d necessitates both the 4-form field strength $\widehat F_4$ and the 4-form potential $\widehat C_4$, whose field strength is the 5-form $\widehat F_5$.
One can write down a ``12d''\footnote{
Not dynamical 12d, but dynamical 10d times a 2-torus.
}
action
\begin{equation}
\begin{aligned}
S_{IIB}=
S_{\text{``12''}}
= \frac{1}{2\kappa_{12}^2}\int_{T_2}dz_1 dz_2\int d^{10}x
\sqrt{-\widehat{g}}
\left(
\widehat{R} - \frac{1}{2}|\widehat F_4|^2-\frac{1}{4}|\widetilde{\widehat{F}}_7|^2
\right)-\frac{1}{8\kappa_{12}^2}\int_{T_2\times \mathbb{R}^{1,9}}\widehat{C}_4\wedge \widehat{F}_4\wedge \widehat{F}_4.
\end{aligned}
\label{12d action}
\end{equation}
The composite seven-form $\widetilde{\widehat{F}}_7$ does not arise from an independent degree of freedom: it is fixed by the self-duality relation \eqref{12d interpretation of 10d self-duality}, $\widetilde{\widehat{F}}_7 = -*_{12}\widetilde{\widehat{F}}_5$. 
For the F1, D1, NS5, D5 and D3 backgrounds treated below, $C_2\wedge H_3$ and $B_2\wedge F_3$ vanish, so $\tilde F_5 = F_5$ and the relation reduces to $\widehat F_7 = -*_{12}d\widehat C_4$.
The action \eqref{12d action} is exactly the Einstein frame type IIB action \eqref{eq:sl2r_invariant_einstein_frame_action}.
The 2d integrand is just a repackaging of $1/\kappa^2$, and it is likely that $\widehat C_3, \widehat C_4$ do not furnish independent degrees of freedom, as has been discussed in \cite{Tseytlin:1996ne}.

\subsection{Uplift of type IIB branes}
We now present the 12d ``uplifts'' of the various type IIB brane solutions, using the effective 12d ``action''.

Although one often thinks of F1 and NS5, D1 and D5 as EM dual pairs, their $H_3$ and $F_3$ configurations are not related by EM duality in 10d, and of course they have opposite dilaton profiles as well.\footnote{The brane solutions are listed in Appendix~\ref{sec:catalog-of-branes}.}

To obtain the NS5 from F1 by EM duality, one actually has to perform an accompanying \SLZ{} transformation to invert the dilaton.
A similar situation is encountered on the D3 brane worldvolume \cite{Tseytlin:1996it}.
Using the unification proposed in the previous subsection, let us uplift the F1, D1, NS5, D5 branes to 12d.

\paragraph{F1}
\begin{equation}
\begin{aligned}
ds^2
&= H^{-3/4}ds_{1,1}^2
+ H^{1/4}ds_8^2
+H^{-1/2}dz_1^2
+H^{1/2}dz_2^2\\
(\widehat{F}_4)_{\mu_0 \mu_1 m z_1}
&= \varepsilon_{\mu_0 \mu_1}\partial_m H^{-1}
\end{aligned}
\end{equation}

\paragraph{D1}
\begin{equation}
\begin{aligned}
ds^2
&= H^{-3/4}ds_{1,1}^2
+ H^{1/4}ds_8^2
+H^{1/2}dz_1^2
+H^{-1/2}dz_2^2\\
(\widehat{F}_4)_{\mu_0 \mu_1 m z_2}
&= \varepsilon_{\mu_0 \mu_1}\partial_m H^{-1}
\end{aligned}
\end{equation}


\paragraph{NS5}
\begin{equation}
\begin{aligned}
ds^2
&= H^{-1/4}ds_{1,5}^2
+ H^{3/4}ds_4^2
+H^{1/2}dz_1^2
+H^{-1/2}dz_2^2\\
(\widehat{F}_4)_{m_1 m_2 m_3 z_1}
&= \varepsilon_{m_1 m_2 m_3 m}\partial_m H
\end{aligned}
\end{equation}

\paragraph{D5}
\begin{equation}
\begin{aligned}
ds^2
&= H^{-1/4}ds_{1,5}^2
+ H^{3/4}ds_4^2
+H^{-1/2}dz_1^2
+H^{1/2}dz_2^2\\
(\widehat{F}_4)_{m_1 m_2 m_3 z_2}
&= \varepsilon_{m_1 m_2 m_3 m}\partial_m H
\end{aligned}
\end{equation}


\paragraph{Comments on role of dilaton and F-theory}
Usually, in supergravity theories, the dilaton is related to the volume of the compact manifold.
In type IIB, however, the compact torus has fixed volume: as the radius of one circle grows, the other shrinks.
The natural place for the non-compact scalar dilaton to live is the radius of one of the circles $R_1 \sim e^{\Phi}$, while the radius of the other circle scales inversely, $R_2 \sim e^{-\Phi}$.
In this way, type IIB backgrounds cannot have a 12d decompactifying limit, because taking the decompactifying limit of one circle immediately compactifies the other circle.
The exception is the D(-1)-instanton uplift, where the pp-wave propagates inside a non-compact torus.
The KK-monopole uplift of the D7-brane is not an exception because a 12d KK-monopole is localised on $S^1 \times \mathbb{R}^3$: there is indeed a compact circle in 12d.

While the two circle radii scale inversely to each other, one may consider the case where they are equal: $e^\Phi = e^{-\Phi} = 1$.
This implies that the radii of both circles are finite, hence a compactification limit of the full torus.
